\documentclass[12pt]{article}
\usepackage{amsmath, amssymb}
\usepackage{geometry}
\usepackage{titlesec}
\usepackage{xcolor}
\usepackage{mdframed}
\usepackage{enumitem}
\usepackage{parskip}
\usepackage{booktabs}

\geometry{margin=1in}

\definecolor{problemblue}{RGB}{30, 80, 160}
\definecolor{solutiongreen}{RGB}{20, 120, 60}
\definecolor{lightblue}{RGB}{230, 240, 255}
\definecolor{lightgreen}{RGB}{230, 248, 235}
\definecolor{boxgray}{RGB}{245,245,245}

\mdfdefinestyle{problembox}{
  backgroundcolor=lightblue,
  linecolor=problemblue,
  linewidth=1.5pt,
  roundcorner=5pt,
  innertopmargin=10pt,
  innerbottommargin=10pt,
  innerleftmargin=12pt,
  innerrightmargin=12pt,
}

\mdfdefinestyle{solutionbox}{
  backgroundcolor=lightgreen,
  linecolor=solutiongreen,
  linewidth=1.5pt,
  roundcorner=5pt,
  innertopmargin=10pt,
  innerbottommargin=10pt,
  innerleftmargin=12pt,
  innerrightmargin=12pt,
}

\mdfdefinestyle{keyresult}{
  backgroundcolor=boxgray,
  linecolor=black,
  linewidth=1pt,
  roundcorner=4pt,
  innertopmargin=8pt,
  innerbottommargin=8pt,
  innerleftmargin=10pt,
  innerrightmargin=10pt,
}

\newcommand{\SE}{[SE]}
\newcommand{\Sg}{[S\gamma]}
\newcommand{\ddt}[1]{\frac{d[#1]}{dt}}
\newcommand{\KME}{K_{M,E}}
\newcommand{\KMg}{K_{M,\gamma}}
\newcommand{\Vmax}{V_{\max}}

\title{
  \textbf{\large CaSB 150 --- Michaelis-Menten Kinetics}\\[6pt]
  \large Problem 4: COVID-19 Spike Protein Reaction System\\[4pt]
  \normalsize Problem Statement + Full Worked Solution
}
\author{}
\date{}

\begin{document}

\maketitle
\vspace{-1em}
\hrule
\vspace{1.5em}

% ============================================================
%  TOPIC OUTLINE
% ============================================================

\section*{Topic Outline: Michaelis-Menten Kinetics}

\subsection*{1. The Reaction Scheme}

The core diagram you must know:
\[
  S + E \;\underset{k_{-1}}{\overset{k_1}{\rightleftharpoons}}\; [SE] \;\overset{k_2}{\longrightarrow}\; P + E
\]
\begin{itemize}[noitemsep]
  \item $S$ = substrate concentration, $E$ = enzyme concentration
  \item $[SE]$ = enzyme-substrate complex (the ``bound state'')
  \item $P$ = product concentration
  \item Three rate constants: $k_1$ (binding), $k_{-1}$ (unbinding), $k_2$ (catalysis / $k_{\text{cat}}$)
  \item \textbf{Key biological fact:} The enzyme is \emph{not consumed} --- it is recycled after each catalytic event.
\end{itemize}

\subsection*{2. Writing the Full System of ODEs}

Four ODEs, one per species:
\begin{align*}
  \frac{d[S]}{dt}  &= -k_1[S][E] + k_{-1}[SE] \\
  \frac{d[E]}{dt}  &= -k_1[S][E] + k_{-1}[SE] + k_2[SE] \\
  \frac{d[SE]}{dt} &= k_1[S][E] - k_{-1}[SE] - k_2[SE] \\
  \frac{d[P]}{dt}  &= k_2[SE]
\end{align*}
Note that $\frac{d[S]}{dt} + \frac{d[SE]}{dt} + \frac{d[P]}{dt} = 0$ (total ``substrate mass'' is conserved).

\subsection*{3. Conservation Laws $\rightarrow$ Reducing the System}

The enzyme is conserved:
\[
  E_0 = [E] + [SE] \quad \Longrightarrow \quad [E] = E_0 - [SE].
\]
Similarly, total substrate mass is conserved: $[S] + [SE] + [P] = \text{const}$.
The product ODE is decoupled (only depends on $[SE]$), so the system reduces to
\textbf{2 independent ODEs} in $[S]$ and $[SE]$.

\subsection*{4. The Quasi-Steady-State Approximation (QSSA)}

\textbf{Key assumption:} $[SE]$ forms and breaks down much faster than $S$ is depleted, so:
\[
  \frac{d[SE]}{dt} \approx 0.
\]
Solving algebraically:
\[
  k_1[S](E_0 - [SE]) - (k_{-1}+k_2)[SE] = 0
  \quad \Longrightarrow \quad
  [SE]^* = \frac{E_0\,[S]}{[S] + K_M},
\]
where the \textbf{Michaelis constant} is:
\[
  K_M = \frac{k_{-1}+k_2}{k_1}.
\]

\subsection*{5. The Michaelis-Menten Rate Law}

The rate of product formation under the QSSA:
\[
  \frac{d[P]}{dt} = k_2[SE]^* = \frac{k_2 E_0\,[S]}{[S]+K_M} = \frac{V_{\max}[S]}{[S]+K_M},
\]
where $V_{\max} = k_2 E_0$ is the maximum production rate (achieved when all enzyme is bound).

\medskip
\textbf{Key limiting behaviors:}
\begin{itemize}[noitemsep]
  \item When $[S] \ll K_M$: rate $\approx \dfrac{V_{\max}}{K_M}[S]$ \quad (linear in $[S]$, first-order kinetics)
  \item When $[S] \gg K_M$: rate $\approx V_{\max}$ \quad (saturated, zero-order kinetics)
  \item When $[S] = K_M$: rate $= V_{\max}/2$ \quad (defines $K_M$ operationally)
\end{itemize}

\subsection*{6. The Hill Function (Generalization)}

When there is \textbf{cooperative binding} (multiple binding sites), the MM rate generalizes to:
\[
  \frac{d[P]}{dt} = \frac{V_{\max}[S]^n}{[S]^n + K_M^n},
\]
where $n$ is the \textbf{Hill coefficient}.
\begin{itemize}[noitemsep]
  \item $n = 1$: recovers standard Michaelis-Menten
  \item $n > 1$: sigmoidal (switch-like) response --- more cooperative binding
  \item $n < 1$: sub-linear saturation
\end{itemize}

\subsection*{7. Extensions to Know for the Exam}

\begin{itemize}
  \item \textbf{Multiple substrates / enzymes / products:} Write one $[SE]$-type ODE per binding pair.
        Apply QSSA to each complex separately. Count independent equations carefully using all
        available conservation laws.
  \item \textbf{Adding degradation/production:} Append $+\lambda$ (production) and $-\delta[\cdot]$
        (first-order degradation) terms to the relevant ODEs. Note that adding production of
        enzyme \emph{breaks} the enzyme conservation law, increasing the number of independent
        equations by one.
\end{itemize}

\subsection*{Exam Skill Checklist (per the 2026 Study Guide)}

\begin{itemize}[noitemsep, label=$\square$]
  \item Write the complete ODE system from a word description or reaction diagram
  \item Identify and apply all conservation laws to reduce the system
  \item Apply the QSSA to solve for $[SE]^*$ and the production rate
  \item Count the number of \textbf{independent} equations (before and after adding degradation/production)
  \item Describe the limiting behaviors and shape of the MM rate curve
  \item Describe the properties of the Hill function and how $n$ changes behavior
\end{itemize}

\vspace{1.5em}
\hrule
\vspace{1.5em}
\newpage

% ============================================================
%  PROBLEM STATEMENT
% ============================================================

\begin{mdframed}[style=problembox]
\textbf{\color{problemblue}\large Problem Statement (20 points)}

\medskip
The COVID-19 RNA vaccine and boosters rely on the spike protein of COVID-19 that is injected
into the body to produce an immune response. Consider the case that spike protein interacts
with a naturally and normally occurring substrate $S$ to produce a product $P_{\text{new}}$
that is not naturally or normally occurring in the human body.
The Michaelis-Menten kinetics for this system can be expressed as:

\medskip
\begin{center}
\renewcommand{\arraystretch}{1.4}
\begin{tabular}{rcccl}
$S + E$        & $\rightleftharpoons$ & $[SE]$      & $\rightarrow$ & $P + E$ \\
$S + \gamma$   & $\rightleftharpoons$ & $[S\gamma]$ & $\rightarrow$ & $P_{\text{new}} + \gamma$
\end{tabular}
\end{center}

\medskip
\noindent where $\gamma$ is the concentration of the spike protein, $P_{\text{new}}$ is the
concentration of the new product, $S$ is the concentration of the substrate, $E$ is the
concentration of the enzyme, $[SE]$ is the concentration of the bound state of the substrate
and enzyme, and $[S\gamma]$ is the concentration of the bound state of the substrate and
spike protein.

\medskip
\begin{enumerate}[label=\textbf{(\alph*)}]
  \item Write down the \textbf{complete} set of equations needed to describe this system in
        terms of Michaelis-Menten kinetics.
  \item How many \textbf{independent} differential equations are needed to describe this
        system?
  \item Solve for the \textbf{Quasi-Steady-State Approximation (QSSA)} for this system.
  \item Over long-time scales, the products will degrade and new substrates and enzymes will
        be created. How would you modify your equations to account for this? How many
        independent equations are now needed?
  \item With the inclusion of degradation of product and production of substrate and enzyme,
        what would you expect the impact of the vaccine to be over short and long times?
        Would it become more or less problematic with time? Which additional biological
        processes might make it more problematic?
\end{enumerate}
\end{mdframed}

\vspace{1.5em}
\hrule
\vspace{1.5em}

% ============================================================
%  BACKGROUND
% ============================================================

\section*{Background: Standard Michaelis-Menten Review}

Before solving the two-pathway problem, recall the single-pathway scheme:
\[
  S + E \;\underset{k_{-1}}{\overset{k_1}{\rightleftharpoons}}\; [SE] \;\overset{k_2}{\longrightarrow}\; P + E.
\]
Here $k_1$ is the binding rate, $k_{-1}$ is the unbinding rate, and $k_2$ (often called
$k_{\text{cat}}$) is the catalytic rate. The Michaelis constant is defined as
\[
  K_M = \frac{k_{-1} + k_2}{k_1},
\]
and the maximum production rate is $V_{\max} = k_2 E_0$, where $E_0$ is the total enzyme
concentration.  In this problem we have \emph{two} such pathways sharing the same substrate
$S$, so we must track both simultaneously.

\vspace{1em}

% ============================================================
%  SOLUTIONS
% ============================================================

\section*{Solutions}

% ---- PART (a) -----------------------------------------------
\subsection*{Part (a): Complete Set of ODEs}

\begin{mdframed}[style=solutionbox]
\textbf{\color{solutiongreen}Strategy.}
Each species appears on the right-hand side of the reaction diagram as either a reactant or
product. We assign rate constants as follows:

\medskip
\begin{center}
\renewcommand{\arraystretch}{1.4}
\begin{tabular}{@{}lll@{}}
\toprule
\textbf{Pathway} & \textbf{Forward (binding)} & \textbf{Reverse / Catalytic} \\
\midrule
Enzyme pathway   & $k_1$ (binding $S+E \to [SE]$)   & $k_{-1}$ (unbinding), $k_2$ (catalysis) \\
Spike pathway    & $k_3$ (binding $S+\gamma \to [S\gamma]$) & $k_{-3}$ (unbinding), $k_4$ (catalysis) \\
\bottomrule
\end{tabular}
\end{center}

\medskip
Now write one ODE for every species, tracking all fluxes in and out.
\end{mdframed}

\medskip
The complete system of ODEs is:

\begin{align}
  \ddt{S} &= -k_1 [S][E] + k_{-1}[SE] - k_3 [S][\gamma] + k_{-3}[S\gamma]
  \label{eq:dS}\\[6pt]
  \ddt{E} &= -k_1 [S][E] + k_{-1}[SE] + k_2[SE]
  \label{eq:dE}\\[6pt]
  \ddt{SE} &= k_1 [S][E] - k_{-1}[SE] - k_2[SE]
  \label{eq:dSE}\\[6pt]
  \ddt{P} &= k_2[SE]
  \label{eq:dP}\\[6pt]
  \ddt{\gamma} &= -k_3 [S][\gamma] + k_{-3}[S\gamma] + k_4[S\gamma]
  \label{eq:dgamma}\\[6pt]
  \ddt{S\gamma} &= k_3[S][\gamma] - k_{-3}[S\gamma] - k_4[S\gamma]
  \label{eq:dSgamma}\\[6pt]
  \frac{d[P_{\text{new}}]}{dt} &= k_4[S\gamma]
  \label{eq:dPnew}
\end{align}

\noindent This is \textbf{7 ODEs} in total describing the dynamics of all 7 species:
$[S],\,[E],\,[SE],\,[P],\,[\gamma],\,[S\gamma],\,[P_{\text{new}}]$.

\medskip

% ---- PART (b) -----------------------------------------------
\subsection*{Part (b): Number of Independent Equations}

\begin{mdframed}[style=solutionbox]
\textbf{\color{solutiongreen}Strategy.}
Count the conservation (constraint) equations that link species together.  Each conservation
law reduces the number of equations by one.
\end{mdframed}

\medskip
There are \textbf{two} conservation laws in this system:

\paragraph{Conservation of total enzyme $E_0$:}
The enzyme $E$ is neither created nor destroyed---it is either free or bound:
\[
  E_0 = [E] + [SE] \quad \Longrightarrow \quad [E] = E_0 - [SE].
\]
This eliminates the ODE for $[E]$.

\paragraph{Conservation of total spike protein $\gamma_0$:}
Similarly, the spike protein is conserved:
\[
  \gamma_0 = [\gamma] + [S\gamma] \quad \Longrightarrow \quad [\gamma] = \gamma_0 - [S\gamma].
\]
This eliminates the ODE for $[\gamma]$.

\paragraph{Elimination of product equations:}
The ODEs for $[P]$ and $[P_{\text{new}}]$ (Eqs.\ \ref{eq:dP} and \ref{eq:dPnew}) are
\emph{decoupled} from the rest of the system---they only appear on the left-hand side. Once
$[SE]$ and $[S\gamma]$ are known, $[P]$ and $[P_{\text{new}}]$ can be obtained by direct
integration. They do not need to be carried as independent dynamical variables.

\paragraph{Summary:}
Starting from 7 ODEs, we eliminate 2 via conservation laws and 2 via decoupling of products,
leaving:

\begin{mdframed}[style=keyresult]
\[
  \boxed{\textbf{3 independent ODEs:} \quad [S],\; [SE],\; [S\gamma].}
\]
\end{mdframed}

These three equations, after substituting the conservation relations, become:
\begin{align}
  \ddt{S} &= -k_1[S](E_0 - [SE]) + k_{-1}[SE] - k_3[S](\gamma_0 - [S\gamma]) + k_{-3}[S\gamma] \\
  \ddt{SE} &= k_1[S](E_0 - [SE]) - (k_{-1} + k_2)[SE] \\
  \ddt{S\gamma} &= k_3[S](\gamma_0 - [S\gamma]) - (k_{-3} + k_4)[S\gamma]
\end{align}

\medskip

% ---- PART (c) -----------------------------------------------
\subsection*{Part (c): Quasi-Steady-State Approximation (QSSA)}

\begin{mdframed}[style=solutionbox]
\textbf{\color{solutiongreen}Strategy.}
The QSSA assumes that the intermediate complexes $[SE]$ and $[S\gamma]$ reach a fast
equilibrium relative to the slow depletion of substrate $S$. We therefore set:
\[
  \frac{d[SE]}{dt} \approx 0 \qquad \text{and} \qquad \frac{d[S\gamma]}{dt} \approx 0,
\]
and solve algebraically for $[SE]^*$ and $[S\gamma]^*$ in terms of $[S]$.
\end{mdframed}

\subsubsection*{Step 1: Solve for $[SE]^*$}

Setting $d[SE]/dt = 0$:
\[
  k_1[S](E_0 - [SE]) - (k_{-1} + k_2)[SE] = 0.
\]
Expand and collect terms in $[SE]$:
\[
  k_1 E_0 [S] = k_1[S][SE] + (k_{-1} + k_2)[SE] = [SE]\bigl(k_1[S] + k_{-1} + k_2\bigr).
\]
Solve:
\[
  [SE]^* = \frac{k_1 E_0 [S]}{k_1[S] + k_{-1} + k_2}.
\]
Dividing numerator and denominator by $k_1$:

\begin{mdframed}[style=keyresult]
\[
  \boxed{[SE]^* = \frac{E_0\,[S]}{[S] + K_{M,E}}}, \qquad K_{M,E} = \frac{k_{-1}+k_2}{k_1}.
\]
\end{mdframed}

\subsubsection*{Step 2: Solve for $[S\gamma]^*$}

By identical algebra, setting $d[S\gamma]/dt = 0$:
\[
  k_3[S](\gamma_0 - [S\gamma]) - (k_{-3} + k_4)[S\gamma] = 0.
\]
\[
  k_3\gamma_0[S] = [S\gamma]\bigl(k_3[S] + k_{-3} + k_4\bigr).
\]

\begin{mdframed}[style=keyresult]
\[
  \boxed{[S\gamma]^* = \frac{\gamma_0\,[S]}{[S] + K_{M,\gamma}}}, \qquad K_{M,\gamma} = \frac{k_{-3}+k_4}{k_3}.
\]
\end{mdframed}

\subsubsection*{Step 3: Production Rates}

The rate of normal product formation and new product formation are:

\begin{align}
  \frac{d[P]}{dt} &= k_2[SE]^* = \frac{k_2 E_0\,[S]}{[S] + \KME} = \frac{V_{\max,E}\,[S]}{[S]+\KME}
  \label{eq:rateP}\\[8pt]
  \frac{d[P_{\text{new}}]}{dt} &= k_4[S\gamma]^* = \frac{k_4\gamma_0\,[S]}{[S]+\KMg} = \frac{V_{\max,\gamma}\,[S]}{[S]+\KMg}
  \label{eq:ratePnew}
\end{align}

where $V_{\max,E} = k_2 E_0$ and $V_{\max,\gamma} = k_4\gamma_0$.

\subsubsection*{Step 4: Effective ODE for $[S]$ under QSSA}

After the QSSA, the substrate $[S]$ is the only remaining dynamical variable:
\begin{equation}
  \ddt{S} = -\frac{V_{\max,E}\,[S]}{[S]+\KME} - \frac{V_{\max,\gamma}\,[S]}{[S]+\KMg}.
  \label{eq:dS_QSSA}
\end{equation}
Both pathways compete to consume $S$, so the total consumption rate is the sum of both MM
terms.

\medskip

% ---- PART (d) -----------------------------------------------
\subsection*{Part (d): Including Degradation and Production}

\begin{mdframed}[style=solutionbox]
\textbf{\color{solutiongreen}Strategy.}
Add first-order degradation terms $(-\delta[\cdot])$ for each product that degrades, and
constant production terms $(+\lambda)$ for substrates and enzymes that are replenished.
\end{mdframed}

Let:
\begin{itemize}[noitemsep]
  \item $\lambda_S$ = production rate of substrate $S$ (units: concentration/time),
  \item $\lambda_E$ = production rate of enzyme $E$,
  \item $\delta_P$ = first-order degradation rate of $P$,
  \item $\delta_{P_{\text{new}}}$ = first-order degradation rate of $P_{\text{new}}$.
\end{itemize}

The modified equations are:
\begin{align}
  \ddt{S}         &= -k_1[S][E] + k_{-1}[SE] - k_3[S][\gamma] + k_{-3}[S\gamma] + \lambda_S \\
  \ddt{E}         &= -k_1[S][E] + k_{-1}[SE] + k_2[SE] + \lambda_E \\
  \ddt{SE}        &= k_1[S][E] - k_{-1}[SE] - k_2[SE] \\
  \ddt{P}         &= k_2[SE] - \delta_P [P] \\
  \ddt{S\gamma}   &= k_3[S][\gamma] - k_{-3}[S\gamma] - k_4[S\gamma] \\
  \frac{d[P_{\text{new}}]}{dt} &= k_4[S\gamma] - \delta_{P_{\text{new}}}[P_{\text{new}}]
\end{align}

\paragraph{Impact on conservation laws.}
Adding $\lambda_E$ to the enzyme ODE \emph{breaks} the conservation law $E_0 = [E]+[SE]$
because $E_0$ is no longer constant---the total enzyme increases over time. Similarly,
$\gamma_0$ remains constant as long as there is no production or degradation of spike protein
itself (the vaccine input is a one-time injection in the simplest model). Therefore:

\begin{mdframed}[style=keyresult]
\[
  \boxed{\textbf{Number of independent equations with degradation/production: 4}}
\]
\[
[S],\;[E] \;\text{(enzyme conservation broken by }\lambda_E),\;[SE],\;[S\gamma].
\]
$[P]$ and $[P_{\text{new}}]$ now couple back in through degradation but can still be tracked
as additional equations if needed (bringing the count to 6 total if products are tracked
explicitly).
\end{mdframed}

\medskip

% ---- PART (e) -----------------------------------------------
\subsection*{Part (e): Biological Interpretation --- Short vs.\ Long Times}

\subsubsection*{Short-time behavior}

At short times after vaccination:
\begin{itemize}
  \item $[\gamma]$ (spike protein) is at its peak concentration---it was just injected.
  \item $[S]$ (natural substrate) is at normal physiological levels.
  \item $[P_{\text{new}}]$ is initially zero but rises rapidly because $[\gamma]$ and $[S]$ are both high.
  \item The production rate of $P_{\text{new}}$ is approximately $\approx V_{\max,\gamma}$ (saturated
        regime) if $[S] \gg \KMg$.
\end{itemize}
\textbf{Conclusion:} The problematic effect ($P_{\text{new}}$ accumulation) is \emph{maximal at short
times}.

\subsubsection*{Long-time behavior}

At long times:
\begin{itemize}
  \item The spike protein $[\gamma]$ is degraded by the immune system, so $[\gamma] \to 0$.
        This drives $[S\gamma] \to 0$ and thus $d[P_{\text{new}}]/dt \to 0$.
  \item Even if $[P_{\text{new}}]$ was produced, it degrades at rate $\delta_{P_{\text{new}}}[P_{\text{new}}]$,
        so $[P_{\text{new}}] \to 0$ as $t \to \infty$.
  \item The substrate $S$ recovers to its steady-state level $\lambda_S / \delta_S$ (if
        degradation of $S$ is also included).
  \item Enzyme $E$ production $\lambda_E$ replenishes any enzyme consumed.
\end{itemize}
\textbf{Conclusion:} The system \emph{becomes less problematic over time} because both the spike
protein and the new product decay.

\subsubsection*{Factors that could make it \emph{more} problematic}

\begin{enumerate}[label=\arabic*.]
  \item \textbf{Slow degradation of $P_{\text{new}}$}: If $\delta_{P_{\text{new}}} \approx 0$, then
        $P_{\text{new}}$ accumulates persistently.
  \item \textbf{High affinity of spike protein for $S$}: A very small $\KMg$ means even tiny
        residual concentrations of $\gamma$ efficiently convert $S$ into $P_{\text{new}}$.
  \item \textbf{Catalytic amplification}: If $P_{\text{new}}$ itself acts as a catalyst for
        additional harmful reactions (a cascade), the damage could grow nonlinearly.
  \item \textbf{Integration of spike-protein genetic material}: If the mRNA were reverse-transcribed
        and integrated into the genome (not expected but sometimes discussed), the body
        could produce spike protein indefinitely, breaking the short-time assumption.
  \item \textbf{Immune suppression of $E$}: If the immune response also degrades or suppresses
        the natural enzyme $E$, the system loses its ability to maintain normal substrate
        homeostasis, potentially causing off-target physiological effects.
\end{enumerate}

\vspace{2em}
\hrule
\vspace{1em}

% ============================================================
%  SUMMARY TABLE
% ============================================================

\section*{Summary of Key Results}

\renewcommand{\arraystretch}{1.5}
\begin{center}
\begin{tabular}{@{}p{3.5cm} p{10cm}@{}}
\toprule
\textbf{Part} & \textbf{Key Result} \\
\midrule
(a) Full ODEs & 7 ODEs for $[S],[E],[SE],[P],[\gamma],[S\gamma],[P_{\text{new}}]$ \\
(b) Independent & \textbf{3 independent ODEs}: $[S],[SE],[S\gamma]$ after applying 2 conservation laws and decoupling products \\
(c) QSSA & $[SE]^* = \dfrac{E_0[S]}{[S]+\KME}$, \quad $[S\gamma]^* = \dfrac{\gamma_0[S]}{[S]+\KMg}$ \\
& Production rates follow standard MM form for each pathway \\
(d) With degradation & \textbf{4 independent ODEs}; enzyme conservation broken by $\lambda_E$ \\
(e) Interpretation & Problematic at short times (high $[\gamma]$); self-resolves at long times as $[\gamma]\to 0$ and $[P_{\text{new}}]$ degrades \\
\bottomrule
\end{tabular}
\end{center}

\end{document}
